
We design experiments to answer:

\textbf{RQ1:} Does fix-impact-ratio discriminate strategic debugging (targeting root causes) from sequential trial-and-error?

\textbf{RQ2:} Do agents cluster errors by shared characteristics (coefficient > 0.3) compared to random ordering?

\textbf{RQ3:} Does error clustering enable prioritization of high-impact fixes (modifications resolving $\geq 2$ test failures)?

\textbf{RQ4:} Do agents transfer learned patterns to held-out test cases at rates exceeding random baseline?

\subsubsection{Datasets}

\textbf{Mock Synthetic Dataset:}
\begin{itemize}
\item \textbf{Problems:} 50 programming tasks (sum, max, array operations, string manipulation, basic algorithms)
\item \textbf{Test Cases:} 15-25 per problem (mean 20.3, median 19)
\item \textbf{Error Types:} Balanced distribution---syntax (24.3\%), runtime (25.3\%), logic (27.4\%), edge\_case (23.0\%)
\item \textbf{Rationale:} Controlled environment validates whether metrics can detect strategic behavior when engineered. Balanced error types ensure clustering signal detectable. Mock validation precedes production deployment.
\end{itemize}

\subsubsection{Baselines}

\textbf{Sequential Trial-and-Error:} Agent addresses test failures one-by-one in test index order, no error clustering or root-cause analysis. Baseline for clustering coefficient (expected $\approx$ 1.0 under random error ordering).

Baselines represent non-strategic approaches, enabling discrimination of clustering and prioritization behaviors.

\subsubsection{Implementation Details}

\textbf{Mock Agent Architecture:}
\begin{itemize}
\item \textbf{Framework:} Python simulation with controlled parameters (clustering\_strength, fix\_success\_rate, pattern\_memory\_enabled)
\item \textbf{Clustering Simulation:} clustering\_strength=0.5 yields 50\% probability of clustering same-type errors consecutively vs random ordering
\item \textbf{Fix Success:} fix\_success\_rate=0.6 (baseline) vs 0.7 (proposed method)
\item \textbf{Iteration Budget:} max\_iterations=10 per problem, temperature=0.7
\item \textbf{Seed:} Fixed random seed (seed=1) for reproducibility
\end{itemize}

\textbf{Hyperparameters (by hypothesis):}
\begin{itemize}
\item \textbf{h-e1:} N=10 problems, controlled trajectories (strategic vs sequential)
\item \textbf{h-m1:} N=50 problems, clustering\_strength=0.5, permutation test with 1000 samples
\item \textbf{h-m2:} N=50 problems, baseline\_fix\_success=0.6, proposed\_fix\_success=0.7, cluster\_bonus\_probability=0.6
\item \textbf{h-m3:} N=50 problems, revealed\_fraction=0.5, pattern\_memory\_enabled=True, max\_iterations=10
\end{itemize}

\textbf{Compute Resources:} Local execution, <1 hour total runtime (mock agents, no GPU required)

\textbf{Reproducibility:} Code available in h-{id}/code/ directories. Implementation uses controlled parameters to validate metric sensitivity, not production agent behavior.

\subsubsection{Evaluation Metrics}

\textbf{Fix-Impact-Ratio:} Tests passed per modification. Strategic agents achieve ratio > 2.0, sequential agents $\approx$ 1.0. Evaluated via t-test or Mann-Whitney U (p < 0.05). Directly measures debugging efficiency---core metric validating strategic vs sequential discrimination.

\textbf{Clustering Coefficient:} Consecutive same-type errors vs expected under random permutation. Coefficient > 0.3 indicates clustering above chance. Evaluated via permutation test (1000 shuffles, p < 0.05). Tests whether agents recognize error patterns, validating mechanism Stage 1.

\textbf{High-Impact Proportion:} Percentage of modifications with $\Delta _passing \geq$ 2 (multi-test fixes). Proposed method expected > baseline. Evaluated via proportion test (p < 0.05). Tests prioritization effectiveness, validating mechanism Stage 2.

\textbf{Held-Out Test Slope Ratio:} Agent slope (held-out pass rate vs iteration) divided by random-mutation slope. Ratio > 1.5 with p < 0.05 indicates transfer learning. Tests pattern transfer to unseen tests without error messages, attempting to validate mechanism Stage 3.

\textbf{Effect Size:} Cohen's d for fix-impact-ratio (0.2=small, 0.5=medium, 0.8=large). Large effect size (d > 0.8) demonstrates strong discriminative power.

Statistical significance threshold p < 0.05 for all tests. Permutation tests control for problem-specific error distributions.
